\section{Evaluation protocol}
\label{sec:protocol}

\subsection{Walk-forward folds}
\label{sec:protocol:folds}
The development period, 2017-02-02 to 2021-12-03, holds \num{1161} trading
days and is divided into four rolling-origin folds. Each fold trains on
\num{500} trading days, selects on the following \num{60}, discards an embargo
of \num{20}, and tests on the next \num{125}. A fifth fold would need
\num{1205} days and is therefore not available. The embargo matters because features use rolling windows of up to
\num{20} days: without it, the input windows of the first test days overlap the
last training days. Test blocks are disjoint by construction, so their return
series concatenate into a single continuous out-of-sample track of \num{500}
trading days.

Rolling rather than anchored windows keep the training sample size constant
across folds, which separates regime effects from sample-size effects.

\subsection{Labels}
\label{sec:protocol:labels}
Within each training date the \num{200} securities with the lowest realised
forward return are labelled $0$ and the \num{200} highest are labelled $1$. The
task is ranking rather than level prediction, and the middle of the
cross-section carries almost no separable signal, so training on the tails gives
a cleaner decision boundary than regression on all names.

\subsection{Training}
\label{sec:protocol:training}
Every condition is trained by the same loop: AdamW (learning rate \num{e-3},
weight decay \num{e-4}) under a cosine schedule over a budget of \num{30}
epochs, batches of \num{512}, gradient norm clipped at \num{1}, binary
cross-entropy on the extreme labels, and a random subsample of \num{60000}
labelled sequences per fold drawn by the seed. Early stopping monitors the
validation AUC with a patience of six epochs and restores the best weights; on
the walk-forward folds it halts after \num{9.9} epochs on average. Because that
could be read as never giving the circuit time to matter, both circuit
conditions are also run with early stopping disabled, for the full
\num{30} epochs (Section~\ref{sec:results:ablation}).

\subsection{Portfolio construction}
Model outputs are standardised within each trading date, restricted to names
that are in the primary universe, not under supervision, and have positive close
and volume, and the top and bottom \num{200} are held long and short with equal
weights. The daily strategy return is the spread between the two legs.

\subsection{Statistics}
\label{sec:protocol:stats}
All Sharpe and Sortino ratios are annualised by $\sqrt{252}$. Information
coefficients are the within-date Spearman correlations between score and
realised forward return; because that series is autocorrelated, its mean is
tested with Newey--West standard errors at the automatic lag
$4(n/100)^{2/9}$~\citep{newey1987simple,newey1994automatic}. Confidence intervals for the Sharpe
ratio come from a stationary bootstrap with mean block length
\num{10}~\citep{politis1994stationary}, which brackets the automatic selection
of \citet{politis2004automatic} on these series; the null is imposed on the
return series by demeaning it before resampling, rather than by re-centring the
bootstrap draws of the statistic, because the latter assumes a symmetric
sampling distribution that a skewed return series does not have.

A non-significant difference is not evidence of no difference, so every null
claim below is accompanied by the two statistics that make it interpretable.
The minimum detectable effect is the smallest true difference the design would
have found at \SI{80}{\percent} power given its standard error. Equivalence is
assessed by two one-sided tests against a stated
margin~\citep{schuirmann1987comparison}, applied to the paired bootstrap
distribution of the Sharpe difference; a significant result there is positive
evidence that the difference lies inside the margin. We pre-specify a margin of
\num{0.03} annualised Sharpe for the trained-versus-frozen contrast, which is
the comparison the design is powered to resolve, and report the wider margins
the other contrasts require rather than claiming equivalence at a margin the
data cannot support.

Two further statistics guard against the selection bias that a study comparing
several model variants necessarily incurs. The Deflated Sharpe
Ratio~\citep{bailey2014deflated} discounts the observed Sharpe by the expected
maximum across the number of configurations searched, which we report
explicitly. The probability of backtest overfitting~\citep{bailey2017probability}
is estimated by combinatorially symmetric cross-validation over the candidate
set; our implementation returns \num{0.53(4)} on independent random strategies,
consistent with the uninformative value of \num{0.5}.

Pairwise contrasts use the Diebold--Mariano test on the difference of daily
return series~\citep{diebold1995comparing}, with the small-sample correction of
\citet{harvey1997testing} because a window of a few hundred days is short for
the normal approximation. That statistic tests a difference in mean return,
which is not the hypothesis the tables are ordered on, so each contrast is also
tested directly on the difference of Sharpe ratios by a paired stationary
bootstrap that resamples both tracks on the same days.

Family-wise error is controlled by the Romano--Wolf stepdown
procedure~\citep{romano2005stepwise}, which is far less conservative than
Bonferroni at this number of comparisons. We state the family explicitly: it is
the set of contrasts against the no-circuit control on one panel, seven
hypotheses. It does not cover the eight tests of each condition against zero,
and it is not pooled across the two markets or the holdout; those are reported
uncorrected and should be read as such.

\subsection{Quantile portfolios}
\label{sec:protocol:quantiles}
The information coefficient is a rank correlation over the whole
cross-section, whereas the long/short book holds only the two tails. The two
can therefore disagree, and on this dataset they do: in the reference
implementation the model with the lowest IC has the highest Sharpe
(Section~\ref{sec:results}). We therefore report decile portfolios --- the mean
forward return of each score decile with Newey--West $t$-statistics, the
top-minus-bottom spread, and the Spearman correlation between decile index and
mean return as a monotonicity measure. A signal with a high IC but low
monotonicity does not belong in a tail-based long/short book, and saying so is
more useful than reporting either statistic alone.

\subsection{Transaction costs}
\label{sec:protocol:costs}
Turnover is the fraction of each leg replaced between consecutive rebalances.
Net returns charge a round-trip cost on the traded fraction of both legs, and we
report Sharpe ratios on a grid of \numlist{0;5;10;20;30;50} basis points
together with the break-even cost at which the mean net return reaches zero.
Reporting gross performance alone is not meaningful for a daily-rebalanced
book: on the reference implementation's own output, turnover averages 24\% per
day and the break-even cost is \SI{7.4}{bps}.

\subsection{Locked holdout}
The supplemental window, 2021-12-06 to 2022-06-24 (\num{135} trading days), is
evaluated exactly once, after every architectural and hyperparameter decision is
frozen. It spans a market regime absent from the development period.
